\documentclass[11pt]{article}
\usepackage{amsmath, amssymb}
\usepackage{booktabs}
\usepackage{graphicx}
\usepackage[margin=1in]{geometry}

\begin{document}

\section*{Problem 4: Programming with CVX}

\subsection*{Part 1}

Using CVXPY, I solved the resource allocation problem with $n = 10$, $D = 10$, and $f_i(x_i) = -ix_i$.

\textbf{Results:}
\begin{itemize}
    \item Optimal value: $-100.0$
    \item Optimal allocation: $x_{10} = 10$, and $x_i = 0$ for $i = 1, \ldots, 9$
\end{itemize}

All resources are allocated to agent 10 because minimizing $\sum_{i} -ix_i$ is equivalent to maximizing $\sum_{i} ix_i$. Since this is a linear program, the optimum lies at a vertex of the feasible polytope, giving everything to the highest-valued agent.

\subsection*{Part 2}

I tested three values of the fairness parameter $\tau$:

\begin{table}[h]
\centering
\begin{tabular}{c|ccc}
\toprule
$i$ & $\tau = 0.1$ & $\tau = 1.0$ & $\tau = 5.0$ \\
\midrule
1  & 0.0111 & 0.1095 & 0.4531 \\
2  & 0.0125 & 0.1229 & 0.4982 \\
3  & 0.0143 & 0.1401 & 0.5533 \\
4  & 0.0166 & 0.1630 & 0.6222 \\
5  & 0.0200 & 0.1947 & 0.7107 \\
6  & 0.0249 & 0.2418 & 0.8284 \\
7  & 0.0332 & 0.3189 & 0.9927 \\
8  & 0.0497 & 0.4681 & 1.2388 \\
9  & 0.0990 & 0.8803 & 1.6468 \\
10 & 9.7187 & 7.3608 & 2.4559 \\
\midrule
\textbf{Optimal Value} & $-95.97$ & $-80.18$ & $-62.96$ \\
\bottomrule
\end{tabular}
\end{table}

\textbf{Observations:}

\begin{itemize}
    \item \textbf{$\tau = 0.1$:} Allocation is nearly identical to Part 1. Agent 10 receives 97.2\% of resources.
    
    \item \textbf{$\tau = 1.0$:} Resources spread more evenly. Agent 10 receives 73.6\%, but all agents get non-trivial allocations.
    
    \item \textbf{$\tau = 5.0$:} Distribution is much more uniform. Agent 10 receives only 24.6\%.
\end{itemize}

The log term promotes fairness because $\log(x) \to -\infty$ as $x \to 0^+$, creating an infinite penalty for giving any agent zero resources. Larger $\tau$ strengthens this fairness effect at the cost of lower utility (optimal value moves from $-100$ toward $-63$).

\end{document}
